\documentclass[10pt,a4paper]{amsart}
\usepackage[margin=19mm]{geometry}
\usepackage{amsmath,amssymb,mathtools}
\usepackage[scr=rsfs,cal=euler]{mathalfa}
\usepackage{xfrac}
\usepackage[unicode,hidelinks,bookmarks=false]{hyperref}
\newcommand{\ie}{i.e.\ }
\newcommand{\cf}{cf.\ }
\newcommand{\resp}{respectively\ }
\newcommand{\Subsection}{\subsection}
\providecommand{\nobreakdash}{\nobreak\hbox{-}}
\setlength{\emergencystretch}{2em}
\multlinegap0pt
\usepackage{amssymb}
\usepackage{tikz}
\newtheorem{lemma}{Lemma}
\title{The Fundamental Limits of Valid Transport Map Estimation}
\author{Sivaraman Balakrishnan}
\date{Source: arXiv:2606.30574v1 (CC BY 4.0)}
\begin{document}
\maketitle

\noindent\emph{Source and licence.} The Fundamental Limits of Valid Transport Map Estimation. Version arXiv:2606.30574v1 (CC BY 4.0), distributed by arXiv under the Creative Commons Attribution 4.0 International licence, http://creativecommons.org/licenses/by/4.0/, as recorded in the arXiv licence metadata for this exact version and shown on the abstract page https://arxiv.org/abs/2606.30574v1. Full licence notice: \texttt{licenses/arxiv-2606.30574-CC-BY-4.0.txt}.\par\smallskip

\noindent\textit{Editorial scope.} Counted unit: the article's lemma lem:lower-bound-analysis (source0.tex lines 1428--1465). The article's complete original proof of it is delivered with it (source0.tex lines 1466--1520, one \texttt{proof} environment of 55 lines). No mathematics is written, repaired or re-derived: the statement and the proof are the article's own lines, in the article's own notation, and no retained line is substituted or re-flowed.  Retained uncounted as the source-local context the proof needs: lines 577--624.  Nothing was omitted from any retained span, so the package carries no omission record.  No retained line contains a citation, so the wrapper carries no bibliography and no bibliography entry.  No retained line contains a figure environment, an image-inclusion command or drawing code, so no image is delivered and the package declares no asset dependency.  Editorial formatting only, in addition: an AMS \texttt{amsart} preamble that restates the article's own declarations and macros used by the retained lines, namely the environments \texttt{lemma} on separate counters, together with the standard packages those lines need.  The retained environments print in this document as Lemma 1 in order of appearance, whereas the article prints the same items with the numbering of its own full document.  The complete native source of the article as distributed in the arXiv e-print for this exact version is bundled unchanged as \texttt{sources/204035/source0.tex}.
\begin{lemma}[Assouad lower bound]
\label{lem:assouad-map-hellinger}
Let $\{P_\sigma:\sigma\in\{\pm1\}^M\}$ be a family of distributions, and let
\[
T_\sigma:\mathcal X\to\mathbb R^d,
\qquad \sigma\in\{\pm1\}^M,
\]
be a corresponding family of maps. Let $\mu$ be a probability measure on
$\mathcal X$, and let $G_1,\ldots,G_M\subseteq\mathcal X$ be disjoint
measurable sets.

For $\sigma\in\{\pm1\}^M$, let $\sigma^{(j)}$ denote the vector obtained
from $\sigma$ by flipping the $j$th coordinate. Suppose that for each
$j=1,\ldots,M$ there exists $\Delta_j>0$ such that
\begin{align}
\label{eq:assouad-local-separation}
\int_{G_j}
\|T_\sigma(x)-T_{\sigma^{(j)}}(x)\|_2^2\,d\mu(x)
\geq
\Delta_j
\end{align}
for every $\sigma\in\{\pm1\}^M$. Suppose also that the Hellinger
affinities satisfy
\begin{align}
\label{eq:assouad-affinity}
\rho(P_\sigma^{\otimes n},P_{\sigma^{(j)}}^{\otimes n})
:=
\int \sqrt{dP_\sigma^{\otimes n} dP^{\otimes n}_{\sigma^{(j)}}}
\geq
\rho_0
\end{align}
for every $\sigma$ and every $j$, where $\rho_0\in(0,1]$.

Then
\begin{align}
\label{eq:assouad-map-conclusion}
\inf_{\widehat T}
\sup_{\sigma\in\{\pm1\}^M}
\mathbb E_\sigma
\int
\|\widehat T(x)-T_\sigma(x)\|_2^2\,d\mu(x)
\geq
\frac{1-\sqrt{1-\rho_0^2}}{8}
\sum_{j=1}^M \Delta_j .
\end{align}
Here the infimum is over all estimators $\widehat T$ based on $n$ 
observations from $P_\sigma$.
\end{lemma}

\begin{lemma}
\label{lem:lower-bound-analysis}
There exists a constant $c_a>0$, depending only on $a$,
such that if
\[
r_n \leq c_a \varepsilon \theta_n,
\]
then, for all sufficiently large $n$, the following statements hold:
\begin{enumerate}

\item For every $\sigma\in\{\pm1\}^n$ and every $j=1,\ldots,n$,
\[
\int_{G_j}
\|T_\sigma(x)-T_{\sigma^{(j)}}(x)\|_2^2\,d\mu(x)
=
4a^2\varepsilon.
\]
Thus the local separation condition in Lemma~\ref{lem:assouad-map-hellinger}
holds with
\[
\Delta_j=4a^2\varepsilon.
\]

\item If $P_\sigma=\nu_\sigma$, then for every $\sigma$ and every $j$,
\[
\rho(P_\sigma^{\otimes n},P_{\sigma^{(j)}}^{\otimes n})
=
(1-\varepsilon)^n.
\]
In particular, since $\varepsilon=1/n$, there is a universal constant
$\rho_0>0$ such that
\[
\rho(P_\sigma^{\otimes n},P_{\sigma^{(j)}}^{\otimes n})
\geq \rho_0
\]
for all sufficiently large $n$.
\end{enumerate}
\end{lemma}

\begin{proof}
Let $\sigma^{(j)}$ be
obtained from $\sigma$ by flipping the $j$th sign. On the left source ball,
\[
y_{j,L}^{\sigma^{(j)}}-y_{j,L}^{\sigma}
=
(0,2\sigma_j a),
\]
and on the right source ball,
\[
y_{j,R}^{\sigma^{(j)}}-y_{j,R}^{\sigma}
=
(0,-2\sigma_j a).
\]
Thus, for every $x\in G_j$,
\[
\|T_\sigma(x)-T_{\sigma^{(j)}}(x)\|_2^2=4a^2.
\]
Since $\mu(G_j)=\varepsilon$, it follows that
\[
\int_{G_j}
\|T_\sigma(x)-T_{\sigma^{(j)}}(x)\|_2^2\,d\mu(x)
=
4a^2\varepsilon.
\]

It remains to verify the Hellinger affinity condition. The distributions
$\nu_\sigma$ and $\nu_{\sigma^{(j)}}$ agree on every block except block
$j$, whose total mass is $\varepsilon$. Inside block $j$, their supports
are disjoint for all sufficiently large $n$, since
$r_n\ll \varepsilon\theta_n\leq \theta_n$ and $a>0$ is fixed. Hence the
one-sample Hellinger affinity is exactly the common mass:
\[
\rho(\nu_\sigma,\nu_{\sigma^{(j)}})
=
1-\varepsilon.
\]
Hellinger affinity tensorizes under products, so
\[
\rho(\nu_\sigma^{\otimes n},\nu_{\sigma^{(j)}}^{\otimes n})
=
\rho(\nu_\sigma,\nu_{\sigma^{(j)}})^n
=
(1-\varepsilon)^n.
\]
Since $\varepsilon=1/n$,
\[
(1-\varepsilon)^n
=
\left(1-\frac1n\right)^n
\]
is bounded below by a positive universal constant for all sufficiently
large $n$. This proves the Hellinger affinity condition and completes the
proof.
\end{proof}

\end{document}
